\scanpage{29}
a. Si $|\lambda|<1$, $\Delta(A)=0$

b. Si $|\lambda|>1$, $\Delta(A)=+\infty$

c. Si $|\lambda|=1$\quad $0\leq\Delta(A)\leq+\infty$, \deleted{$\Delta(A)<\infty$}, et $\Delta(A)<+\infty\Longleftrightarrow B\in\underline a$.

\deleted{Ces \unsure{résultats} peu pratiques, la considération de}

Par suite, $0\leq\Delta(1+A)<+\infty$ si $A\in\underline a$, \deleted{et}
\deleted{\unsure{pratique seule la considération} de $\Delta:\ 1+\underline a$}
mais on \unsure{considérera}, $\Delta(A)$ comme définie \deleted{\unsure{à une indét.}} \unsure{lorsqu'}il est \deleted{de la forme $1+B$, $B\in\underline a$} \inserted{\unsure{contenu} dans $1+\underline a$}.

\uuline{Prop. 8}
\diagram{\begin{tikzcd}[column sep=large,row sep=small]
\left\{\begin{aligned}
\Delta(A)&=\Delta(A^*)\\
\Delta(AB)&=\Delta(A)\Delta(B)\\
\Delta(1)&=1.
\end{aligned}\right. & {}\\
{}&\inserted{A,B\in1+\underline a}\arrow[ul,bend left=20,end anchor=north east]
\end{tikzcd}}

\deleted{Si $A\in\mathcal A$, soit $\unsure{\sigma(A)|r(A)}$ le projecteur orthogonal sur $A\mathcal H$, on a $\unsure{\sigma(A)|r(A)}\in\mathcal A$. \deleted{\unsure{Supposons}} \deleted{\unsure{fini}} \deleted{$A\in\underline a^+$, alors}
\[
(5)\qquad\Delta(1+A)=\Delta_E(1_E+A_E)\qquad E=\unsure{\sigma|s}.
\]
où $A_E$ et $1_E$ désignent $E$ et $A$ \unsure{considérés} comme éléments de l'algèbre de von Neumann $\mathcal A_E=E\mathcal A E$, et $\Delta_E$ la fonction \unsure{correspondante} dans $\mathcal A_E$ muni de la trace induite par celle de $\mathcal A$.}

Soit $E\in P(\mathcal A)$, posons $\mathcal A_E=E\mathcal A E$ \inserted{sous-algèbre de $\mathcal A$, ayant \unsure{l'unité} $E$, et \unsure{alors}}, c'est une algèbre de von Neumann dans $E.\mathcal H$, \deleted{\unsure{dont}} on la munit de la trace \inserted{$T_E$} induite par $T$, on écrira $\Delta_{\deleted{E}}^{\inserted{E}}(A)$ pour le $\Delta$ d'un $A\in\mathcal A_E$ \inserted{\unsure{qui est défini}} \unsure{considéré sur} $\mathcal A$
\mspage{30}
